% English draft based on appendices/appendix_b_homology_ja.tex
\section{Cross-Domain Homology}

The paper repeatedly uses homology rather than identity. The claim is not that physical particles, organisms, conscious states, intelligent stopping layers, and social institutions are the same thing. The claim is that they instantiate a repeated generative pattern.

\[
\text{incomplete closure}
\;\longrightarrow\;
\text{emergent attractor}
\;\longrightarrow\;
\text{sustainable dynamics}.
\]

\begin{longtable}{@{}p{0.20\linewidth}p{0.34\linewidth}p{0.36\linewidth}@{}}
\toprule
Layer & Non-closure & Attractor / stopping form \\
\midrule
\endfirsthead
\toprule
Layer & Non-closure & Attractor / stopping form \\
\midrule
\endhead
Physics & differential flow and quasi-closure & particle-like attractor \\
Life & non-equilibrium informational exchange & life attractor $\Omega_{\text{life}}$ \\
Consciousness & temporal non-closure and speed-scale difference & conscious state as windowed stabilization \\
Intelligence & no internal fixed point in $L_F$ & emergent stopping layer $L_E$ \\
Society & distributed disagreement and non-final legitimacy & $\Phi_{\text{net}}$ synchronization \\
\bottomrule
\end{longtable}

Homology preserves differences. It identifies the repeated pattern without reducing one layer to another.

\bigskip
